%% 参赛信息表：封面下部内容（由 biomed-report-theme.sty 的 \makereporttitle 装配）。
%% 原“参赛作品简介”章的表格，2026-08 移到封面；无章标题、无目录项。
{\renewcommand{\arraystretch}{1.6}%
\begin{tabular}{|>{\centering\arraybackslash}m{0.18\textwidth}|m{0.76\textwidth}|}
\hline
\textbf{参赛作品名称} & BioMed QAgent：面向生物医学研究的多源数据智能采集与分析系统 \\ \hline
\textbf{参赛选题} & 赛道2-方向1A \\ \hline
\textbf{参赛作品简介} & 系统采用“开放式推理与确定性数据处理分权”设计：Qwen 等 LLM 驱动的 Pi Agent 负责理解研究目标、发现候选来源并提交声明式数据需求；TypeScript Dataset Core 负责来源登记、解析执行、规范化、多源整合、质量门禁与正式发布。系统不把 Agent 工作区的临时 CSV 视为任务完成，而以含数据表、Schema、Provenance、Audit、Validation 与内容摘要的不可变 Publication 作为正式结果，并以动态 Family 协议表达多表拓扑。来源经 SourceAsset、SHA-256 与 SourceLocator 进入证据链；不确定映射、未知单位与低置信度抽取可触发人在回路请求；中断后依据事件与 checkpoint 恢复，无需模型重新解释已确认的变换。 \\ \hline
\textbf{参赛作品使用Qwen大模型及AI技术说明} & 本作品未将 Qwen 写死在配置中，而是提供 API key 配置方法，允许自选含 Qwen key 在内的 LLM 作主 Agent。实验（第五章）中，Qwen 系列承担 gold8 正式发布（qwen3.8-27B）、图表理解（DashScope Qwen-VL）与全谱系对抗验证（qwen3.7/3.8 各档位）使用 DashScope key 时可调用 Qwen-VL 理解页面图像进行数据提取。 \\ \hline
\textbf{作品相关文件链接} & \url{https://github.com/Lidozs55/BioMed-QAgent} \\ \hline
\end{tabular}}
